\section{Finite arithmetic and combinatorics}
\label{sec:finite}

\subsection{A primitive recursive definition}

Write $[n]=\{0,\ldots,n-1\}$. A \emph{restricted growth word} of length $n$ is a word
$a_0\cdots a_{n-1}$ with $a_0=0$ when $n>0$ and
\[
  a_i\leq 1+\max_{j<i}a_j\qquad(0<i<n).
\]
The empty word is allowed. Its number of blocks is zero; that of a nonempty
word is $1+\max_i a_i$. Encode a word by its $n$ base-$(n+1)$ digits, retaining
the length separately. Validity and block count are bounded-loop tests, and
there are at most $(n+1)^n$ candidate codes. Let $S(n,k)$ count valid words
with $k$ blocks, and put
\begin{equation}\label{eq:bell-stirling}
  B_n=\sum_{k=0}^{n}S(n,k).
\end{equation}
In particular $B_0=1$. This definition is primitive recursive without any
unbounded search.

The blocks of a word are the nonempty sets $\{i:a_i=j\}$. Conversely, order the
blocks of a partition by their least elements and give them labels
$0,1,\ldots$. These two algorithms are inverse. Consequently $B_n$ counts
partitions of $[n]$, equivalence relations on $[n]$, and rhyme schemes after
canonical relabelling. An equivalence relation is encoded by an $n$ by $n$
Boolean matrix; its classes are found by scanning for their least members.
This also proves the equivalence-relation interpretation without forming a
quotient of an infinite set.

\begin{proposition}[Stirling recurrence and positivity]\label{prop:stirling}
With $S(n,k)=0$ outside $0\leq k\leq n$,
\[
 S(0,0)=1,\qquad S(n+1,k)=kS(n,k)+S(n,k-1).
\]
Thus $B_n\geq1$, and, for $n\geq1$,
\[
 2^{n-1}\leq B_n\leq n^n.
\]
The first eleven values are
\[
 1,1,2,5,15,52,203,877,4140,21147,115975.
\]
\end{proposition}
\begin{proof}
The last element of a partition either joins one of its $k$ pre-existing
blocks or is the sole member of a new block. Deleting that element gives the
displayed recurrence and its inverse construction. The all-singleton
partition proves positivity. For the lower bound, choose any subset of the
other $n-1$ elements to join the block containing $0$, and make every remaining
element a singleton. These $2^{n-1}$ partitions are distinct. Every restricted
growth word of positive length $n$ is a word over $[n]$, giving the upper bound.
The numerical values follow by the recurrence, using only integer addition
and multiplication.
\end{proof}

If $q_1,\ldots,q_n$ are distinct primes, partitions of $[n]$ correspond to
unordered factorizations of $q_1\cdots q_n$ into factors exceeding one: a block
is sent to the product of the primes it contains. In the inverse algorithm,
test divisibility of every factor by each of the given primes. The elementary
prime-divisor lemma, applied finitely often, makes these supports disjoint and
exhaustive. Thus the squarefree-factorization statement also has a finite
proof. The input includes the distinct primes; no general factorization oracle
is needed.

\subsection{Recurrences and finite summation formulas}

The two familiar recurrences are unified by Spivey's finite identity
\cite{spivey}; the proof below gives the explicit bijection.

\begin{theorem}[Bell recurrence and Spivey's identity]\label{thm:spivey}
For $n,m\geq0$, using $0^0=1$,
\begin{align}
 B_{n+1}&=\sum_{k=0}^{n}\binom nk B_k,\label{eq:bell-binomial}\\
 B_{n+m}&=\sum_{k=0}^{n}\sum_{j=0}^{m}
       \binom nk S(m,j)j^{n-k}B_k.\label{eq:spivey}
\end{align}
\end{theorem}
\begin{proof}
For the first identity, delete the entire block containing a distinguished
element of an $(n+1)$-element set. If $k$ elements remain, their set can be
chosen in $\binom nk$ ways and partitioned in $B_k$ ways. Putting the
distinguished element together with the complement reverses the operation.

For the second identity, first restrict a partition to a fixed $m$-element
subset. If this restriction has $j$ blocks, there are $S(m,j)$ choices for it.
Among the other $n$ elements choose the $k$ elements whose blocks do not meet
the first subset. Partition those elements in $B_k$ ways. Each of the remaining
$n-k$ elements chooses one of the $j$ old blocks. These choices are independent,
and give $j^{n-k}$ possibilities. Restricting a completed partition recovers
all the choices uniquely. The convention $0^0=1$ handles empty subsets.
\end{proof}

Here and below a combinatorial count is an equality between bounded sums of
Boolean characteristic functions. Each displayed inverse map gives a
primitive recursive permutation of the relevant finite codes; equality of
their sums follows by finite reindexing.

\begin{proposition}[Bell triangle]\label{prop:bell-triangle}
Define $T(0,0)=1$ and, for $0\leq k\leq n$,
\[
  T(n,0)=B_n,\qquad
  T(n,k)=T(n,k-1)+T(n-1,k-1)\quad(k>0).
\]
Then
\[
  T(n,k)=\sum_{i=0}^{k}\binom ki B_{n-i},\qquad
  T(n,n)=B_{n+1}.
\]
In particular the next row begins with the preceding row's last entry.
\end{proposition}
\begin{proof}
Pascal's identity shows that the displayed sum satisfies the defining
addition rule, including its initial column. At $k=n$, replace $i$ by $n-i$
and use \eqref{eq:bell-binomial}. These verifications use induction through a
finite triangular table.
\end{proof}

For an indeterminate $X$, put $(X)_0=1$ and
$(X)_k=X(X-1)\cdots(X-k+1)$. Comparing the recurrence in
Proposition~\ref{prop:stirling} with
$X(X)_k=(X)_{k+1}+k(X)_k$ proves the polynomial identity
\begin{equation}\label{eq:power-falling}
 X^n=\sum_{k=0}^{n}S(n,k)(X)_k.
\end{equation}
All polynomial identities here are statements about finite coefficient
lists, not assertions requiring a theory of arbitrary functions.

\begin{theorem}[Inclusion--exclusion and finite Dobinski forms]
\label{thm:finite-dobinski}
For all $n,k\geq0$,
\begin{equation}\label{eq:stirling-ie}
 k!S(n,k)=\sum_{i=0}^{k}(-1)^{k-i}\binom ki i^n.
\end{equation}
Consequently
\begin{align}
 B_n&=\sum_{k=0}^{n}\frac1{k!}
        \sum_{i=0}^{k}(-1)^{k-i}\binom ki i^n,\label{eq:bell-ie}\\
 B_n&=\sum_{i=0}^{n}\frac{i^n}{i!}
        \sum_{j=0}^{n-i}\frac{(-1)^j}{j!}.\label{eq:bell-finite-dobinski}
\end{align}
If $D_t$ is the number of fixed-point-free permutations of $[t]$, then
\begin{equation}\label{eq:bell-derangements}
  B_n=\sum_{i=0}^{n}\frac{i^nD_{n-i}}{i!(n-i)!}.
\end{equation}
\end{theorem}
\begin{proof}
There are $k!S(n,k)$ surjections from $[n]$ to $[k]$: label the blocks of a
partition bijectively by the target labels. For a fixed function with image
$I$, its contribution to the right side of \eqref{eq:stirling-ie} is
\[
 \sum_{I\subseteq A\subseteq[k]}(-1)^{k-|A|}
       =(1-1)^{k-|I|}.
\]
This is one exactly when the function is onto and zero otherwise. Interchanging
the two finite sums proves \eqref{eq:stirling-ie}; summing over $k$ proves
\eqref{eq:bell-ie}. Put $j=k-i$ and use
$\binom ki/k!=1/(i!j!)$ to obtain
\eqref{eq:bell-finite-dobinski}.

The same inclusion--exclusion calculation for forbidden fixed points gives
$D_t=t!\sum_{j=0}^{t}(-1)^j/j!$: imposing a specified set of $j$ fixed points
leaves $(t-j)!$ permutations. Substitution gives
\eqref{eq:bell-derangements}. These rational formulas are finite arithmetic
equalities; multiplication by a product of the relevant factorials removes
all division.
\end{proof}

One may also group partitions by their block sizes. If $m_j$ blocks have size
$j$, then the labels can be assigned to ordered slots in $n!$ ways. Permuting
within each block and permuting equal-sized blocks gives exactly
$\prod_{j=1}^{n}(j!)^{m_j}m_j!$ assignments with the same partition. Hence
\begin{equation}\label{eq:block-profile}
 B_n=\sum_{\substack{m_1,\ldots,m_n\leq n\\\sum_{j=1}^n jm_j=n}}
       \frac{n!}{\prod_{j=1}^{n}(j!)^{m_j}m_j!}.
\end{equation}
The quotient is exact, justified by this explicit finite grouping. More
generally, weighting a block of size $j$ by $x_j$ defines the complete Bell
polynomial
\[
 C_n(x_1,\ldots,x_n)=\sum_{\pi}\prod_{A\in\pi}x_{|A|}.
\]
Its coefficient sum is $C_n(1,\ldots,1)=B_n$.
The one-variable Touchard polynomial
$\Phi_n(t)=\sum_kS(n,k)t^k$ similarly satisfies
\[
 \Phi_n(1)=B_n,\qquad
 \Phi_{n+1}(t)=t\bigl(\Phi_n(t)+\Phi_n'(t)\bigr),
\]
by the Stirling recurrence. The derivative is the formal derivative of a
finite polynomial.

\subsection{A finite moment functional and the shift identities}
\label{subsec:finite-functional}

Because the falling-factorial polynomials are monic of successively increasing
degrees, every polynomial $f\in\Q[X]$ has a unique finite expansion
$f=\sum_jc_j(X)_j$. Define
\[
 \mathcal L(f)=\sum_jc_j.
\]
Triangular coefficient elimination computes $\mathcal L$. Identity
\eqref{eq:power-falling} says $\mathcal L(X^n)=B_n$, while by definition
$\mathcal L((X)_j)=1$.

\begin{lemma}[Finite shift identity]\label{lem:finite-shift}
For every polynomial $f$ and every $k\geq0$,
\begin{equation}\label{eq:finite-shift}
 \mathcal L\bigl((X)_k f(X)\bigr)
          =\mathcal L\bigl(f(X+k)\bigr).
\end{equation}
\end{lemma}
\begin{proof}
First let $g(Y)=(Y)_j$. Then
$(X)_k g(X-k)=(X)_{k+j}$, so applying $\mathcal L$ gives one on both sides of
$\mathcal L((X)_k g(X-k))=\mathcal L(g(X))$. Expand an arbitrary $g$ in the
falling-factorial basis and use linearity. Finally set $g(Y)=f(Y+k)$.
\end{proof}

This single finite lemma proves the additional summation identities on the
Wikipedia page, attributed there to Komatsu--Pita-Ruiz~\cite{komatsu-pita}.
For rational $a,b$ and natural $k,n$,
\begin{equation}\label{eq:general-umbral-shift}
 \mathcal L\bigl((aX+b)^n\bigr)
 =\mathcal L\bigl((X)_k(aX+b-ak)^n\bigr).
\end{equation}
Define the signed first-kind Stirling coefficients by
$(X)_k=\sum_i s(k,i)X^i$, equivalently
$s(k,i)=(-1)^{k-i}\genfrac{[}{]}{0pt}{}{k}{i}$. Expanding both sides gives
\begin{align}\label{eq:general-stirling-shift}
 &\sum_{j=0}^{n}\binom nj a^j b^{n-j}B_j\notag\\
 &\quad=\sum_{i=0}^{k}s(k,i)
       \sum_{j=0}^{n}\binom nj a^j(b-ak)^{n-j}B_{j+i}.
\end{align}
In particular $k=1$ gives the displayed one-step identity of the source,
and $a=1,b=k$ gives
\[
 \sum_{j=0}^{n}\binom nj k^{n-j}B_j
       =\sum_{i=0}^{k}s(k,i)B_{n+i}.
\]
Using \eqref{eq:finite-shift} with $k=1$ and
$f(X)=X^n(X-1)^r$ also gives
\[
 \mathcal L\bigl(X^{n+1}(X-1)^r\bigr)
       =\mathcal L\bigl(X^r(X+1)^n\bigr).
\]
On expanding, this is
\begin{equation}\label{eq:general-pascal-inversion}
 \sum_{j=0}^{n}\binom nj B_{r+j}
   =\sum_{i=0}^{r}(-1)^{r-i}\binom ri B_{n+i+1}.
\end{equation}
Its specialization $n=0$ is Pascal inversion for the Bell recurrence.
No infinite moment representation was needed for these identities.

\subsection{The Hankel determinant and strict log-convexity}

\begin{theorem}[Hankel determinant]\label{thm:hankel}
For every $N\geq0$,
\begin{equation}\label{eq:hankel}
 \det\bigl(B_{i+j}\bigr)_{0\leq i,j\leq N}
       =\prod_{j=0}^{N}j!.
\end{equation}
\end{theorem}
\begin{proof}
Define monic Charlier polynomials by the finite formula
\[
 P_r(X)=\sum_{j=0}^{r}(-1)^{r-j}\binom rj(X)_j.
\]
Let $\Delta f(X)=f(X+1)-f(X)$. The finite shift identity and the binomial
formula for an iterated difference yield
\[
 \mathcal L(P_r f)
   =\sum_{j=0}^{r}(-1)^{r-j}\binom rj\mathcal L(f(X+j))
   =\mathcal L(\Delta^r f).
\]
If $\deg f<r$, its $r$th difference is zero. If $f$ is monic of degree $r$,
the $r$th difference is the constant $r!$. Applying this to the $P_s$, and
using symmetry when $r<s$, gives
\[
  \mathcal L(P_rP_s)=\begin{cases}r!&r=s,\\0&r\ne s.\end{cases}
\]
The coefficient matrix $U$ expressing $P_0,\ldots,P_N$ in the monomial basis
is lower triangular with diagonal entries one. Thus, by multiplication of
finite matrices,
\[
 U\,(B_{i+j})_{i,j=0}^{N}\,U^{\mathsf T}
       =\operatorname{diag}(0!,1!,\ldots,N!).
\]
Taking determinants, defined by their finite permutation sums, proves
\eqref{eq:hankel}, since $\det U=1$.
\end{proof}

\begin{corollary}[Strict log-convexity]\label{cor:strict-logconvex}
For $n\geq1$,
\begin{equation}\label{eq:strict-logconvex}
 B_n^2<B_{n-1}B_{n+1}.
\end{equation}
Consequently the quotients $B_{n+1}/B_n$ increase, and
$B_mB_n\leq B_{m+n}$ for $m,n\geq0$.
\end{corollary}
\begin{proof}
If $f=\sum_jc_jP_j$, the orthogonality calculation gives
$\mathcal L(f^2)=\sum_jj!c_j^2$, positive for nonzero $f$.
For two independent polynomials $f,g$, write their coefficient vectors in this
basis as $c,d$. Then the entirely finite identity
\[
 \mathcal L(f^2)\mathcal L(g^2)-\mathcal L(fg)^2
       =\sum_{i<j}i!j!(c_id_j-c_jd_i)^2>0
\]
follows by expansion. At least one minor is nonzero, since otherwise the
vectors would be proportional. The same conclusion holds for the bilinear
form $\mathcal L(Xfg)$, because \eqref{eq:finite-shift} identifies it with
$\mathcal L(f(X+1)g(X+1))$, and substitution $X\mapsto X+1$ preserves
independence.

Write $n-1=2q+\epsilon$ with $\epsilon\in\{0,1\}$. Take $f=X^q$ and
$g=X^{q+1}$, using the ordinary bilinear form for $\epsilon=0$ and the shifted
one for $\epsilon=1$. Their Gram determinant is precisely
$B_{n-1}B_{n+1}-B_n^2$, proving strict positivity. Finally multiply the
inequalities
$B_{n+j+1}/B_{n+j}\geq B_{j+1}/B_j$ for $0\leq j<m$; the products telescope.
All denominators are positive integers and can be cleared.
\end{proof}

\subsection{Normalized log-concavity, with a finite proof}
\label{subsec:logconcavity}

The complementary inequality also admits a finite proof. We give the
coefficient argument underlying the Bender--Canfield method
\cite{bender-canfield}, specialized enough that its arithmetic content is
visible.

\begin{lemma}[A finite coefficient inequality]\label{lem:coefficient-logconcave}
Let $x_0=1,x_1,x_2,\ldots$ be nonnegative rational numbers such that
\[
 x_jx_{k-1}\geq x_{j-1}x_k\qquad(1\leq j<k).
\]
Let $a_0=1$, put $a_t=0$ for $t<0$, and define
\[
 r a_r=\sum_{j=1}^{r}x_j a_{r-j}\qquad(r\geq1).
\]
Then, for $1\leq m\leq n$,
\begin{equation}\label{eq:general-logconcavity}
 a_m a_n\geq a_{m-1}a_{n+1}.
\end{equation}
Only $x_0,\ldots,x_{n+1}$ and $a_0,\ldots,a_{n+1}$ occur in each assertion.
\end{lemma}
\begin{proof}
The recurrence, with its $x_0a_r$ term restored, says
$\sum_{j=1}^{r+1}x_{j-1}a_{r+1-j}=(r+1)a_r$.
Expand the product of this identity and the defining recurrence, subtract
the two resulting products, and pair indices $j,k$. One obtains the finite
identity
\begin{align}\label{eq:paired-logconcavity}
 &m(n+1)(a_m a_n-a_{m-1}a_{n+1})\notag\\
 &\quad=\sum_{1\leq j<k\leq n+1}
 (x_jx_{k-1}-x_{j-1}x_k)
 (a_{m-j}a_{n+1-k}-a_{m-k}a_{n+1-j}).
\end{align}
Indeed, before pairing, the summand is
$(x_jx_{k-1}-x_{j-1}x_k)a_{m-j}a_{n+1-k}$; summands with $j>m$ vanish by the
negative-index convention, and exchanging $j,k$ negates the first factor.

Induct on $m$, retaining all $n\geq m$ as free parameters. If $k>m$, the
second subtracted product in \eqref{eq:paired-logconcavity} is zero, so its
bracket is nonnegative. Suppose $k\leq m$, and put
\[
 A=m-j,\quad C=m-k,\quad D=n+1-j,\quad E=n+1-k.
\]
Then $C\leq\min(A,E)$, $\max(A,E)\leq D$, and $A+E=C+D$.
Starting with the pair $(\min(A,E),\max(A,E))$, move one unit at a time
from its smaller index to its larger one until the pair becomes $(C,D)$.
At each step the smaller index is less than $m$, so the induction hypothesis
says the product of the two $a$'s cannot increase. Thus
$a_Aa_E\geq a_Ca_D$. The case $m=1$ has only the already-settled
$k>m$ terms. Every summand in \eqref{eq:paired-logconcavity} is therefore
nonnegative. Since $m(n+1)>0$, this proves the lemma.

For a literal quantifier-free induction implementation, fix a final bound
$N$ and retain the Boolean conjunction of these inequalities for all
$1\leq m\leq n\leq N$ reached so far. The above steps never require an index
exceeding $N+1$. Clearing the product of denominators makes the invariant an
integer inequality test on a finite table.
\end{proof}

\begin{theorem}[Normalized log-concavity]\label{thm:normalized-logconcavity}
For every $n\geq1$,
\begin{equation}\label{eq:normalized-logconcavity}
 n B_{n-1}B_{n+1}\leq(n+1)B_n^2.
\end{equation}
Equivalently, $B_n/n!$ is log-concave. Moreover,
\begin{equation}\label{eq:bell-product-bounds}
 B_mB_n\leq B_{m+n}\leq\binom{m+n}{m}B_mB_n.
\end{equation}
\end{theorem}
\begin{proof}
Set $a_n=B_n/n!$, $x_0=1$, and $x_j=1/(j-1)!$ for $j\geq1$.
Dividing \eqref{eq:bell-binomial} by $n!$ gives
$(n+1)a_{n+1}=\sum_{j=1}^{n+1}x_j a_{n+1-j}$.
The successive ratios $x_j/x_{j-1}$ are
$1,1,1/2,1/3,\ldots$, a decreasing sequence, so the hypothesis of
Lemma~\ref{lem:coefficient-logconcave} holds. Its case $m=n$ is exactly
\eqref{eq:normalized-logconcavity} after factorials are cleared.

For the upper product bound, normalized log-concavity makes $a_{j+1}/a_j$
decrease. Hence
\[
 \frac{a_{n+m}}{a_n}
 =\prod_{j=0}^{m-1}\frac{a_{n+j+1}}{a_{n+j}}
 \leq\prod_{j=0}^{m-1}\frac{a_{j+1}}{a_j}=a_m.
\]
Clearing factorials gives the upper inequality in
\eqref{eq:bell-product-bounds}; the lower one was proved in
Corollary~\ref{cor:strict-logconvex}. The empty product handles $m=0$.
\end{proof}

\subsection{Permutation interpretations and card shuffles}

A generalized pattern $1$--$23$ means indices $i<j<j+1$ with
$\pi_i<\pi_j<\pi_{j+1}$; $1$--$32$ instead means
$\pi_i<\pi_{j+1}<\pi_j$. These definitions are bounded tests on a permutation.

\begin{proposition}[Eight avoidance classes]\label{prop:pattern-avoidance}
For each of
\[
 1\text{--}23,\quad32\text{--}1,\quad3\text{--}21,\quad12\text{--}3,
 \qquad
 1\text{--}32,\quad23\text{--}1,\quad3\text{--}12,\quad21\text{--}3,
\]
the number of avoiding permutations of $[n]$ is $B_n$.
\end{proposition}
\begin{proof}
Given a partition, order its blocks by decreasing minima. For the first
class, write each block with its minimum first and its remaining elements in
decreasing order, then erase the block separators. An ascent can occur only
immediately after a left-to-right minimum, so it cannot be preceded by a
smaller element. The result avoids $1$--$23$. Conversely, cut such an avoiding
permutation immediately before each left-to-right minimum. Within a segment,
every adjacent pair after its first entry must decrease: otherwise the first
entry would complete a forbidden pattern. The segments therefore recover
exactly the original partition and its prescribed order.

For the class $1$--$32$, use the same block order but write every block in
increasing order. A descent can then end only at a left-to-right minimum.
This is exactly the condition excluding $1$--$32$: a descent ending at an
element that is not a new minimum has an earlier smaller element and gives
the forbidden pattern. Cutting at the minima again supplies the inverse.
Reversing the positions and complementing every value $v$ to $n-1-v$ are
explicit involutions on permutations. Their four combinations take the first
pattern to the first four listed patterns, and the second to the last four.
\end{proof}

\begin{proposition}[Extending every $321$ to a $3241$]
\label{prop:barred-pattern}
There are $B_n$ permutations of $[n]$ in which every triple of positions
$i<j<k$ with $\pi_i>\pi_j>\pi_k$ has a position $j<\ell<k$ satisfying
$\pi_\ell>\pi_i$.
\end{proposition}
\begin{proof}
Cut a permutation immediately before each left-to-right maximum. Call these
its record segments; each begins with its largest element. The stated
property is equivalent to the tail of every record segment being increasing.
For necessity, a decreasing pair in such a tail, together with the leading
record, is a $321$ that cannot be extended: no larger value occurs before
the end of the segment. Conversely, consider any $321$. Its middle element
is not a left-to-right maximum. If its last element lay in the same segment,
the increasing-tail condition would be violated. Otherwise a new record lies
between the middle and last elements, and this record exceeds the first
element of the triple, supplying the required fourth position.

Given a set partition, order blocks by increasing maxima and write each block
with its maximum first and the remaining elements in increasing order.
Erasing separators produces exactly a permutation with the just-described
record segments. Cutting at the records is its inverse. This is the required
finite bijection.
\end{proof}

\begin{proposition}[Return shuffles]\label{prop:return-shuffles}
For a deck of $n\geq1$ distinct cards, there are $n^n$ sequences of $n$
operations that remove the top card and reinsert it in any of the $n$
positions. Exactly $B_n$ return the deck to its original order. Thus the
return probability under independent uniform choices is $B_n/n^n$.
\end{proposition}
\begin{proof}
Reverse a path of operations. An inverse operation chooses a card and moves
it to the top. Label the original order $1,\ldots,n$. After a word of such
choices, the chosen cards occur first, in reverse order of their last
occurrences in the word; untouched cards follow in their original order.
This description follows by induction on the length of the word.

To return to $1,\ldots,n$, the chosen labels must therefore be a prefix
$\{1,\ldots,k\}$, and their last occurrences must be ordered from latest to
earliest as $1,\ldots,k$. Partition the $n$ action times according to equal
chosen labels. Conversely, given any partition of those times, label its
blocks $1,\ldots,k$ in decreasing order of their maxima. The resulting choice
word returns the deck to its original order. The constructions are inverse;
reversing operations now proves the assertion for top-card reinsertion.
\end{proof}

The permutation bijection gives $B_n\leq n!$. A second useful comparison is
$n^n\leq n!B_n$: partition functions $[n]\to[n]$ by their fibers and use
\eqref{eq:power-falling} at $X=n$, noting $(n)_k\leq n!$.
For $n\geq3$ this is strict, since the $k=1$ term has
$(n)_1=n<n!$. Consequently the return probability exceeds $1/n!$ for
$n\geq3$; equality holds for $n=1,2$.

\paragraph{Arithmetic audit.}
The largest finite objects used in this section are tables, words,
permutations, and coefficient lists with explicitly bounded lengths.
Their enumerations have primitive recursive bounds, such as $(n+1)^n$, $n!$,
and products of factorials. Every appeal to induction can retain a finite
table of the preceding assertions as one quantifier-free invariant. The
positive-denominator convention turns all rational inequalities into integer
inequalities. Thus the arguments use no unrestricted induction, choice,
excluded middle for undecidable formulas, or existence principle for infinite
sets.

% Bibliography metadata:
% bender-canfield: E. A. Bender and E. R. Canfield,
% Log-Concavity and Related Properties of the Cycle Index Polynomials,
% J. Combin. Theory Ser. A 74 (1996), 57--70,
% https://mathweb.ucsd.edu/~ebender/reprints/93.pdf
% Additional finite-sum statements match T. Komatsu and C. Pita-Ruiz,
% Some formulas for Bell numbers, Filomat 32(11) (2018), 3881--3889,
% https://doi.org/10.2298/FIL1811881K . Their proofs above follow directly
% from the finite shift functional, without using the article as a premise.
